\section{Hướng phát triển và giai đoạn thực nghiệm tiếp theo}
\label{sec:future}

Kết quả v4.1/v5.1 xác định hai vấn đề ưu tiên: path ranking, đặc biệt đối với
reverse reasoning, và độ nhạy thấp của benchmark đối với structural
intervention. Giai đoạn tiếp theo cần cải thiện hai vấn đề này một cách tách
biệt để tránh gộp thay đổi retrieval với thay đổi verification engine.

\subsection{Query-aware retrieval và path reranking}

Answer intent và vai trò nguyên nhân--trung gian--hệ quả nên được xác định
trước khi chọn primary path. Đối với reverse reasoning, outcome được phát
hiện trong truy vấn cần được cố định làm điểm cuối, sau đó hệ thống thực hiện
backward search có ràng buộc để tìm predecessor. Path reranker có thể kết hợp
semantic score với các đặc trưng lexical và cấu trúc, bao gồm số điều, mức
hình phạt, vị trí event trên path, rule condition/effect và độ đặc hiệu của
outcome.

Mỗi thay đổi retrieval cần được đánh giá trước verification trên cùng
candidate pool. Các metric chính gồm Rule/Event Recall, Oracle Path, Top-1
Path và ranking error. Chỉ sau khi retrieval được đóng băng mới chạy lại
paired structural SCM--path ablation, qua đó giữ được khả năng quy kết chênh
lệch cho verification engine.

\subsection{Benchmark nhạy với intervention}

Một benchmark structural-counterfactual mới nên biểu diễn mỗi mẫu dưới dạng:

\begin{equation}
\left(
q,\,
C_F,\,
\doop(M=m),\,
Y_F,\,
Y_{\doop(M=m)},\,
z
\right),
\label{eq:future-counterfactual-item}
\end{equation}

trong đó $q$ là truy vấn, $C_F$ là factual context, $M$ là mediator bị can
thiệp, $Y_F$ và $Y_{\doop(M=m)}$ lần lượt là factual và counterfactual
outcome, còn $z$ là nhãn quan hệ. Tối thiểu, benchmark cần cân bằng ba nhãn:

\begin{itemize}
    \item \texttt{NECESSARY}: can thiệp lên mediator làm thay đổi outcome;
    \item \texttt{NON\_NECESSARY}: outcome vẫn được duy trì bởi một mechanism
    thay thế đang hoạt động;
    \item \texttt{INDETERMINATE}: dữ kiện, root state hoặc rule không đủ để
    kết luận.
\end{itemize}

Các nhóm \texttt{SUFFICIENT}, \texttt{NO\_EFFECT} và conflict có thể được
bổ sung để kiểm tra nhiều nhánh của engine. Benchmark phải chứa các cặp mà
node deletion và hard intervention đưa ra kết quả khác nhau; nếu không,
paired ablation vẫn không phân biệt được semantics của hai phương pháp.

Mỗi mẫu cần được chuyên gia kiểm tra ở cả factual world và intervened world.
Gold annotation nên bao gồm rule được dùng, rule bị vô hiệu hóa, mechanism
thay thế, state của mediator/outcome và giải thích bằng ngôn ngữ tự nhiên.

\subsection{Mở rộng biểu diễn quy tắc và LegalSCM}

Mô hình tiếp theo cần hỗ trợ điều kiện hội, phủ định, ngoại lệ, rule priority
và hiệu lực theo thời gian. Một rule tổng quát có thể được biểu diễn bằng
nhiều literal ở vế điều kiện thay vì một positive condition duy nhất. Khi có
xung đột, hệ thống cần phân biệt rõ \texttt{FALSE} với
\texttt{UNKNOWN}, đồng thời ghi lại quy tắc nào tạo ra từng state.

Event schema cũng cần tách canonical ID, alias và preferred display name.
Alias được dùng cho retrieval, còn preferred name được dùng trong câu trả
lời. Thay đổi này vừa giảm nhiễu semantic vừa tránh đưa chuỗi
\texttt{||} vào output.

\subsection{Đánh giá cấu trúc phi tuyến và citation}

Branch và convergence nên được đánh giá bằng path-set hoặc subgraph metric,
không ép thành một chuỗi duy nhất. Các metric dự kiến gồm branch outcome
coverage, convergence input coverage, subgraph edge F1 và graph edit
distance.

Đối với citation, cần chuyển từ so khớp tập article sang assertion-level
grounding. Mỗi assertion trong short answer hoặc explanation phải được ánh
xạ tới rule và article hỗ trợ. Chuyên gia sau đó có thể đánh giá ba tiêu chí
riêng biệt: article có liên quan, article có entail assertion và citation có
đủ để hỗ trợ kết luận hay không.

\subsection{Đánh giá bên ngoài và đánh giá con người}

Thực nghiệm tương lai cần bổ sung các baseline được chạy dưới cùng corpus và
giao thức, bao gồm vanilla RAG, graph RAG không có verification, LLM-only và
các phương pháp causal/counterfactual RAG có thể tái lập. Một test set độc
lập do chuyên gia biên soạn cần được tách khỏi graph-construction pipeline.

Ngoài metric tự động, nghiên cứu cần đo liệu LegalSCM trace có thực sự giúp
người đánh giá:

\begin{itemize}
    \item phát hiện path hoặc citation sai nhanh hơn;
    \item hiểu sự khác nhau giữa factual và counterfactual world;
    \item nhận biết mechanism thay thế;
    \item hiệu chỉnh mức tin cậy vào câu trả lời.
\end{itemize}

\subsection{Tăng cường khả năng tái lập}

Trước một lần chạy kết luận, cần cố định commit code, revision của embedding
model, dependency lockfile, random seed, cấu hình server và hash của toàn bộ
input/output artifact. Hai mode phải bắt đầu từ cùng cached retrieval output
và ghi paired difference ở mức từng sample. Runtime nên được lặp nhiều lần
trong điều kiện cold cache và warm cache.

Các giả thuyết, primary metric, exclusion rule và subgroup analysis cần được
xác định trước khi xem kết quả cuối. Cách tổ chức này giúp ngăn việc lựa chọn
metric hậu nghiệm và làm rõ liệu cải thiện đến từ retrieval, verification hay
answer generation.